\newpage
\section{Fluxos do Sistema}

\tikzset{
    c4b/.style={draw, rounded corners=3pt, minimum width=2.3cm, minimum height=0.95cm, align=center, font=\footnotesize, fill=white, line width=0.6pt},
    c4a/.style={-{Latex[length=2mm]}, line width=0.6pt},
    c4d/.style={-{Latex[length=2mm]}, line width=0.6pt, dashed},
    c4x/.style={-{Latex[length=2mm]}, line width=0.6pt, dash pattern=on 2pt off 1.5pt on 0.6pt off 1.5pt},
    c4l/.style={line width=0.6pt},
    c4n/.style={fill=white, inner sep=1.2pt, font=\scriptsize},
}
\newcommand{\cNav}{\node[c4b] (nav) at (0,0) {Navegador};}
\newcommand{\cWeb}{\node[c4b] (web) at (3.6,0) {Front-end};}
\newcommand{\cApi}{\node[c4b] (api) at (7.2,0) {Back-end};}
\newcommand{\cWrk}{\node[c4b] (wrk) at (10.8,0) {\emph{Worker}};}
\newcommand{\cRts}{\node[c4b] (rts) at (14.4,0) {Executor\\de Scripts};}
\newcommand{\cPg}{\node[c4b, minimum width=5cm] (pg) at (9,2.6) {PostgreSQL};}
\newcommand{\cRed}{\node[c4b, minimum width=5cm] (red) at (9,-3.2) {Redis};}
\newcommand{\cAcq}{\node[c4b] (acq) at (9,-5.8) {Aquisição};}
\newcommand{\cCam}{\node[c4b] (cam) at (0,-5.8) {Campo\\(CLP)};}
\newcommand{\cs}[5][0.5]{\draw[#2] (#3) -- node[c4n, pos=#1] {#5} (#4);}
\newcommand{\cVt}[6]{\draw[#1] ([xshift=#4]#2) -- node[c4n, pos=#6] {#5} ([xshift=#4]#2 |- #3);}
\newcommand{\cVtt}[8]{\draw[#1] ([xshift=#4]#2) -- node[c4n, pos=#6] {#5} node[c4n, pos=#8] {#7} ([xshift=#4]#2 |- #3);}
\newcommand{\cHz}[6]{\draw[#1] ([yshift=#4]#2) -- node[c4n, pos=#6] {#5} ([yshift=#4]#2 -| #3);}
\newcommand{\cHzz}[8]{\draw[#1] ([yshift=#4]#2) -- node[c4n, pos=#6] {#5} node[c4n, pos=#8] {#7} ([yshift=#4]#2 -| #3);}
\newcommand{\cUp}[7]{\coordinate (cuA) at ([xshift=#4]#2.north); \coordinate (cuB) at ([xshift=#5]#3.north); \draw[#1] (cuA) -- ($(cuA)+(0,#6)$) coordinate (cuC) -- node[c4n] {#7} (cuC -| cuB) -- (cuB);}
\newcommand{\cDn}[7]{\coordinate (cuA) at ([xshift=#4]#2.south); \coordinate (cuB) at ([xshift=#5]#3.south); \draw[#1] (cuA) -- ($(cuA)-(0,#6)$) coordinate (cuC) -- node[c4n] {#7} (cuC -| cuB) -- (cuB);}

\par Este capítulo descreve os caminhos que o dado e as ações percorrem no sistema: quem pede, quem responde, onde cada informação é guardada e o que acontece quando uma peça falha. Cada fluxo é apresentado por um diagrama de blocos com setas numeradas, seguido de um texto que percorre os passos na ordem.

\par Os blocos correspondem aos serviços da arquitetura e têm as posições fixas da Figura~\ref{fig:fx_mapa}. O navegador e o equipamento de campo ficam fora do servidor. Dentro dele estão o \emph{front-end}, que entrega as telas, o \emph{back-end}, que concentra as regras de funcionamento, o \emph{worker}, que retira tarefas das filas e grava no banco, a aquisição, que conversa com o equipamento, e o executor de scripts, que roda, de forma isolada, o JavaScript escrito pelos usuários. O PostgreSQL guarda os dados que precisam durar. O Redis guarda o valor atual de cada tag, as filas e as travas, e também transporta a comunicação entre os serviços.

\begin{figure}[H]
    \centering
    \caption{Mapa dos blocos usados nos fluxos}
    \label{fig:fx_mapa}
    \vspace{0.4cm}
    \figpadrao{\begin{tikzpicture}
        \cNav \cWeb \cApi \cWrk \cRts \cRed \cPg \cAcq \cCam
        \begin{scope}[on background layer]
            \node[draw, dashed, rounded corners=6pt, fit=(web)(api)(wrk)(rts)(red)(pg)(acq), inner sep=9pt, fill=black!4] (srv) {};
        \end{scope}
        \node[font=\scriptsize, anchor=south east] at (srv.south east) {Servidor};
        \draw[c4l] (nav) -- (web);
        \draw[c4l] (nav.north) -- ++(0,0.8) -| ([xshift=0.5cm]api.north west);
        \draw[c4l] (api.south) -- (api.south |- red.north);
        \draw[c4l] (wrk.south) -- (wrk.south |- red.north);
        \draw[c4l] (api.north) -- (api.north |- pg.south);
        \draw[c4l] (wrk.north) -- (wrk.north |- pg.south);
        \draw[c4l] (wrk) -- (rts);
        \draw[c4l] (red) -- (acq);
        \draw[c4l] (acq) -- (cam);
    \end{tikzpicture}}
    \vspace{0.2cm} \\
    Fonte: Autor.
\end{figure}

\par Em todas as figuras, a linha cheia indica o caminho normal, a tracejada indica uma resposta e a de traço e ponto indica um caminho alternativo, percorrido quando alguma peça falha. Os números entre parênteses seguem a ordem dos passos, e o mesmo número em dois trechos de linha indica um único passo que atravessa o Redis, como nas chamadas entre o \emph{back-end} e a aquisição, que o Moleculer transporta por esse banco \cite{MOLECULER_DOCS,REDIS_DOCS}. Os passos alternativos recebem o número do passo normal e uma letra, como em (4b). Cada fluxo termina com o que acontece quando uma peça deixa de funcionar.

\subsection{Cadastro de uma Tag}

\par O cadastro de uma tag é o ponto em que o sistema passa a conhecer uma nova variável do processo. A tag pertence a um dispositivo, ligado a uma planta, e traz o mapeamento do ponto no equipamento e as regras de histórico. O cadastro é feito com o sistema em funcionamento, e a nova tag chega à aquisição sem reiniciar nenhum processo, como mostra a Figura~\ref{fig:fx_cadastro_tag}.

\begin{figure}[H]
    \centering
    \caption{Fluxo do cadastro de uma tag}
    \label{fig:fx_cadastro_tag}
    \vspace{0.4cm}
    \figpadrao{\begin{tikzpicture}
        \cNav \cWeb \cApi \cRed \cPg \cAcq
        \cs{c4a}{nav}{web}{(1)}
        \cUp{c4a}{nav}{api}{0.4cm}{-0.6cm}{0.55}{(2)}
        \cUp{c4d}{api}{nav}{0.1cm}{-0.4cm}{1.1}{(4)}
        \cVt{c4a}{api.south}{red.north}{0cm}{(5)}{0.5}
        \cVt{c4a}{api.north}{pg.south}{0.8cm}{(3)}{0.5}
        \cVt{c4a}{red.south}{acq.north}{-0.35cm}{(6)}{0.3}
        \cVt{c4a}{acq.north}{red.south}{0.35cm}{(7)}{0.7}
    \end{tikzpicture}}
    \vspace{0.2cm} \\
    Fonte: Autor.
\end{figure}

\par O administrador abre a tela de tags de um dispositivo, entregue pelo \emph{front-end} (1). Ao salvar o formulário, o navegador envia o pedido ao \emph{back-end} (2), que confere a permissão do usuário e as regras do cadastro e grava a tag e o registro da mudança no banco (3).

\par O \emph{back-end} responde ao navegador (4) e publica o evento de mudança do cadastro (5), que o Moleculer entrega, pelo Redis, ao Agendador, dentro da aquisição (6). Como a aquisição não tem banco de dados, o Agendador pede ao \emph{back-end} a configuração pronta do dispositivo (7) e inclui a tag no plano de leitura. A primeira leitura acontece no ciclo seguinte, e até lá a tela mostra a tag como sem dado.

\par Quando a validação falha, o \emph{back-end} recusa o pedido com uma mensagem, e nada é gravado. Se o processo cai entre a gravação e a publicação do evento, a mudança continua registrada e é publicada depois.

\par Se o \emph{back-end}, o Redis ou o banco param, o cadastro fica indisponível, mas a leitura do campo continua com a última configuração recebida. Se a aquisição para, o cadastro continua funcionando, e ela recebe as mudanças quando volta.

\subsection{Do Campo ao Banco}

\par Este é o fluxo principal de um supervisório, da leitura no equipamento até o histórico no banco, como mostra a Figura~\ref{fig:fx_campo}.

\begin{figure}[H]
    \centering
    \caption{Fluxo da leitura do equipamento até o banco}
    \label{fig:fx_campo}
    \vspace{0.4cm}
    \figpadrao{\begin{tikzpicture}
        \cWrk \cRed \cPg \cAcq \cCam
        \cHz{c4a}{acq.west}{cam.east}{0.25cm}{(1)}{0.5}
        \cHz{c4d}{cam.east}{acq.west}{-0.25cm}{(2)}{0.5}
        \cDn{c4x}{cam}{acq}{0cm}{0cm}{0.8}{(2b)}
        \cVt{c4a}{acq.north}{red.south}{0cm}{(3)}{0.5}
        \cVt{c4a}{red.north}{wrk.south}{1.4cm}{(4)}{0.5}
        \cVt{c4a}{wrk.south}{red.north}{0.6cm}{(5)}{0.5}
        \cVt{c4a}{wrk.north}{pg.south}{0cm}{(6)}{0.5}
    \end{tikzpicture}}
    \vspace{0.2cm} \\
    Fonte: Autor.
\end{figure}

\par No período de cada dispositivo, a aquisição envia uma requisição Modbus ao \sigla{clp} (1) \cite{MODBUS_TCP_IMPLEMENTATION_GUIDE}, que responde com os valores brutos (2). A aquisição registra o instante da resposta, atribui a qualidade e publica o resultado em duas filas no Redis, uma de valor atual, chamada de ingestão, e outra de histórico (3).

\par O \emph{worker} retira as leituras das filas (4), grava o valor atual no Redis (5) e grava o histórico no banco (6), conforme o modo da tag: não gravar, gravar toda leitura ou gravar quando o valor muda além da banda morta.

\par Se o equipamento não responde ou devolve um erro (2b), a aquisição publica as tags com qualidade ruim, e o \emph{worker} mantém o último valor, com a indicação da falha. Se as falhas se repetem, o evento de falha de comunicação do dispositivo é ativado, e o primeiro ciclo bem-sucedido o encerra.

\par Se o \emph{worker} para, as leituras se acumulam nas filas e são processadas quando ele volta. Se o Redis volta vazio, uma rotina refaz o valor atual a partir do banco. Se o banco fica indisponível, os trabalhos de histórico esperam e tentam de novo.

\subsection{Do Banco à Tela em Tempo Real}

\par A tela obtém o valor atual por consultas periódicas: ela pergunta os valores das tags da planta aberta, e a resposta vem apenas do Redis, sem consultar o banco nem o equipamento, como resume a Figura~\ref{fig:fx_polling}.

\begin{figure}[H]
    \centering
    \caption{Fluxo da exibição do valor atual por consulta periódica}
    \label{fig:fx_polling}
    \vspace{0.4cm}
    \figpadrao{\begin{tikzpicture}
        \cNav \cApi \cRed
        \cHz{c4a}{nav.east}{api.west}{0.25cm}{(1)}{0.5}
        \cHz{c4d}{api.west}{nav.east}{-0.25cm}{(3)}{0.5}
        \cVt{c4a}{api.south}{red.north}{0cm}{(2)}{0.5}
    \end{tikzpicture}}
    \vspace{0.2cm} \\
    Fonte: Autor.
\end{figure}

\par A tela envia o pedido ao \emph{back-end} (1), que lê do Redis o valor atual de cada tag (2) e define sua qualidade. Se nada mudou, a resposta é vazia. Caso contrário, o \emph{back-end} devolve os valores (3), e a tela atualiza o que mudou.

\par A qualidade aparece com texto e ícone. O valor parado há mais de três períodos aparece como desatualizado, e a tag cuja leitura falha mostra o último valor bom com um símbolo de alerta e o texto de falha de comunicação. A lista de acompanhamento, o sinótico, o medidor do painel e a grade de tags usam a mesma consulta e a mesma forma de exibição. Depois de três falhas seguidas de pedido, a tela mostra o aviso de conexão instável e aumenta o intervalo entre as tentativas.

\par Com o Redis fora do ar ou vazio, ou com o banco fora do ar, a tela mostra que o serviço está indisponível em vez de valores. A queda do equipamento de campo não derruba a tela, que continua a responder com o último valor e a qualidade correspondente.

\subsection{Histórico e Tendência}

\par A consulta ao histórico alimenta os gráficos de tendência e permite ver períodos anteriores, como mostra a Figura~\ref{fig:fx_tendencia}. A tela envia o pedido de tendência (1), com as tags e o intervalo desejado. O \emph{back-end} consulta o banco (2) e devolve o resultado (3), e a tela desenha a série. Depois da primeira carga, a tela pede apenas o trecho novo. Um pedido grande demais é recusado, e a falha da consulta não prejudica a exibição dos valores atuais, que têm caminho próprio.

\begin{figure}[H]
    \centering
    \caption{Fluxo da consulta de tendência}
    \label{fig:fx_tendencia}
    \vspace{0.4cm}
    \figpadrao{\begin{tikzpicture}
        \cNav \cApi \cPg
        \cHz{c4a}{nav.east}{api.west}{0.25cm}{(1)}{0.5}
        \cHz{c4d}{api.west}{nav.east}{-0.25cm}{(3)}{0.5}
        \cVt{c4a}{api.north}{pg.south}{0cm}{(2)}{0.5}
    \end{tikzpicture}}
    \vspace{0.2cm} \\
    Fonte: Autor.
\end{figure}

\subsection{Comando de Equipamentos}

\par O comando é a escrita de um valor em um ponto do equipamento, como ligar um motor ou alterar um ponto de ajuste. Na tela, o operador escolhe a tag, informa o novo valor e confirma, e o navegador envia o pedido ao \emph{back-end} (1), como mostra a Figura~\ref{fig:fx_comando_execucao}.

\begin{figure}[H]
    \centering
    \caption{Fluxo do comando}
    \label{fig:fx_comando_execucao}
    \vspace{0.4cm}
    \figpadrao{\begin{tikzpicture}
        \cNav \cApi \cRed \cPg \cAcq \cCam
        \cHz{c4a}{nav.east}{api.west}{0.25cm}{(1)}{0.5}
        \cHz{c4d}{api.west}{nav.east}{-0.25cm}{(10)}{0.5}
        \cVtt{c4a}{api.north}{pg.south}{0.4cm}{(2)}{0.3}{(9)}{0.7}
        \cVtt{c4a}{api.south}{red.north}{-0.3cm}{(3)}{0.3}{(4)}{0.7}
        \cVt{c4x}{api.south}{red.north}{0.3cm}{(8b)}{0.4}
        \cVt{c4d}{red.north}{api.south}{-0.9cm}{(8)}{0.35}
        \cVt{c4a}{red.south}{acq.north}{-0.6cm}{(4)}{0.5}
        \cVt{c4a}{acq.north}{red.south}{0cm}{(5)}{0.3}
        \cVt{c4d}{acq.north}{red.south}{0.6cm}{(8)}{0.7}
        \cHz{c4a}{acq.west}{cam.east}{0.25cm}{(6)}{0.5}
        \cHz{c4d}{cam.east}{acq.west}{-0.25cm}{(7)}{0.5}
        \cDn{c4x}{cam}{acq}{0cm}{0cm}{0.8}{(7b)}
    \end{tikzpicture}}
    \vspace{0.2cm} \\
    Fonte: Autor.
\end{figure}

\par O \emph{back-end} confere a permissão do usuário, se a tag permite escrita, se o dispositivo se comunica e se o valor está dentro da faixa, e registra o comando no banco (2). Sem a certeza de que o equipamento se comunica, o comando é recusado. Um pedido repetido não cria um segundo comando, e cada tag tem no máximo um comando em andamento.

\par O \emph{back-end} obtém no Redis a trava da tag (3) e chama a ação de escrita da aquisição pelo Moleculer (4), sem fila e sem nova tentativa automática. A aquisição obtém no Redis a trava do dispositivo (5), pois o comando tem prioridade sobre a leitura periódica, e envia a escrita ao \sigla{clp} (6). O sucesso é a resposta normal do equipamento, que repete o pedido enviado (7) e é conferida contra o endereço e o valor, e a próxima leitura periódica mostra o efeito real. A aquisição devolve o resultado ao \emph{back-end} (8), que fecha o comando, grava o registro (9) e informa o diálogo (10). O registro guarda quem emitiu o comando, o valor, o horário e o resultado, e não pode ser apagado nem alterado.

\par Além do usuário, um script também pode escrever, inclusive em hora marcada, por um evento agendado. Essa escrita usa o mesmo caminho do comando e cede ao comando do operador. A tag interna, que não tem equipamento, recebe o valor diretamente, sem passar pela aquisição.

\par Se o equipamento responde com exceção, com uma repetição que não confere ou não responde no prazo (7b), a escrita é tratada como falha, e a mensagem pede que o operador confira o valor do ponto. Se o sistema perde o rastro, por exemplo, se a resposta da aquisição não chega ao \emph{back-end} (8b), o comando é fechado como sem confirmação do equipamento, e a tag fica bloqueada por trinta e cinco segundos. Quando uma peça falha, o sistema prefere recusar a arriscar: sem o Redis, o Agendador ou a aquisição, o comando é recusado.

\subsection{Alarme}

\par O alarme transforma uma condição do processo em um aviso que exige a atenção de uma pessoa. A condição é um evento da tag: a tag, um operador de comparação, o valor de referência e uma banda morta opcional, que evita que um valor oscilando em torno do limite abra e feche o alarme. O alarme liga um ou mais eventos a um de quatro níveis fixos, Crítico, Alto, Médio e Baixo, como mostra a Figura~\ref{fig:fx_alarme}.

\begin{figure}[H]
    \centering
    \caption{Fluxo da detecção, da sinalização e do reconhecimento de um alarme}
    \label{fig:fx_alarme}
    \vspace{0.4cm}
    \figpadrao{\begin{tikzpicture}
        \cNav \cApi \cWrk \cRed \cPg \cAcq
        \cVt{c4a}{acq.north}{red.south}{0cm}{(1)}{0.5}
        \cVt{c4a}{red.north}{wrk.south}{2.1cm}{(2)}{0.5}
        \cVtt{c4a}{wrk.south}{red.north}{-0.3cm}{(3)}{0.3}{(5)}{0.7}
        \cVt{c4x}{wrk.south}{red.north}{-0.9cm}{(4b)}{0.5}
        \cVt{c4a}{wrk.north}{pg.south}{0cm}{(4)}{0.5}
        \cHzz{c4a}{nav.east}{api.west}{0cm}{(6)}{0.3}{(8)}{0.7}
        \cVt{c4a}{api.south}{red.north}{0cm}{(7)}{0.5}
        \cVt{c4a}{api.north}{pg.south}{0cm}{(9)}{0.5}
    \end{tikzpicture}}
    \vspace{0.2cm} \\
    Fonte: Autor.
\end{figure}

\par A leitura é publicada pela aquisição na fila de ingestão (1) e retirada pelo \emph{worker} (2), que, depois de gravar o valor atual, avalia os eventos da tag. A avaliação é feita no Redis, que guarda o estado de cada detector (3): compara o valor com o limite, aplica a banda morta e decide se houve transição. Havendo transição, o \emph{worker} a grava no banco (4), abrindo, atualizando ou fechando a ocorrência do alarme, e atualiza no Redis o sumário de alarmes (5).

\par A moldura da aplicação, presente em todas as páginas, pergunta periodicamente o sumário de alarmes (6), que o \emph{back-end} lê no Redis (7). O indicador exibe o nível mais grave não reconhecido e a quantidade de ocorrências. O operador reconhece uma ocorrência na página de alarmes (8), e o \emph{back-end} grava o reconhecimento, com o nome da pessoa (9).

\par O evento composto combina outros eventos com os operadores E, OU e NÃO, e o evento agendado dispara em horários definidos. Os eventos de falha de comunicação, de leitura e de escrita são criados automaticamente com o dispositivo Modbus. Se o processo cai antes da gravação no banco (4b), a transição não se perde: uma tarefa periódica grava o que ficou pendente no Redis.

\par Com o Redis fora, a avaliação para, e o trabalho é repetido. Com o banco indisponível, as transições ficam pendentes até ele voltar. Com o \emph{worker} parado, nenhum alarme novo é detectado, e a fila preserva as leituras.

\subsection{Variável Meta}

\par A variável meta é uma tag cujo valor vem de um cálculo escrito pelo administrador sobre outras tags da mesma planta, como uma diferença de pressões. O código, em JavaScript, roda no executor de scripts, isolado, sem banco, rede, arquivo nem segredo, como mostra a Figura~\ref{fig:fx_meta}.

\begin{figure}[H]
    \centering
    \caption{Fluxo do cálculo de uma variável meta}
    \label{fig:fx_meta}
    \vspace{0.4cm}
    \figpadrao{\begin{tikzpicture}
        \cWrk \cRed \cPg \cAcq
        \node[c4b] (rts) at (15.2,0) {Executor\\de Scripts};
        \cVt{c4a}{acq.north}{red.south}{0cm}{(1)}{0.5}
        \cVt{c4a}{red.north}{wrk.south}{2.1cm}{(2)}{0.5}
        \cVtt{c4a}{wrk.south}{red.north}{-0.3cm}{(3)}{0.3}{(7)}{0.7}
        \cVt{c4a}{wrk.north}{pg.south}{0cm}{(4)}{0.5}
        \cHz{c4a}{wrk.east}{rts.west}{0.34cm}{(5)}{0.5}
        \cHz{c4d}{rts.west}{wrk.east}{0cm}{(6)}{0.5}
        \cHz{c4x}{rts.west}{wrk.east}{-0.34cm}{(6b)}{0.5}
    \end{tikzpicture}}
    \vspace{0.2cm} \\
    Fonte: Autor.
\end{figure}

\par A variável é recalculada a cada leitura nova de uma tag que o código lê ou em horários definidos. A aquisição publica a leitura na fila de ingestão (1), e o \emph{worker} a retira (2), identifica quais variáveis meta leem aquela tag e coloca os cálculos em uma fila própria (3).

\par O \emph{worker} retira os cálculos e busca os dados de que o código precisa: o valor atual, no Redis, e os pontos de histórico, no banco (4). Os dados seguem com o código ao executor de scripts, pelo Moleculer (5). Para cada cálculo, o executor cria uma cópia nova do interpretador, com limites de tempo, memória e saída, e devolve o valor ou o erro (6).

\par O \emph{worker} valida o que veio e o publica como leitura da tag interna (7), pelo mesmo caminho das leituras do equipamento. A qualidade da meta é a pior entre as tags que ela lê. O mesmo executor roda a tag interna gerada, o script de evento e o componente de conteúdo do sinótico. Se o código ultrapassa um limite ou gera um erro (6b), apenas aquela meta recebe erro, e o \emph{worker} publica a qualidade ruim em vez de um valor.

\par Se o executor cai, nenhuma meta é calculada, e a variável fica com o último valor. As demais tags, seus alarmes e os comandos não dependem dele.

\subsection{Efeito das Falhas}

\par A Tabela~\ref{tab:fx_falhas} reúne o efeito da falha de cada bloco e mostra que o sistema perde apenas as funções ligadas à peça que falhou.

{\renewcommand{\arraystretch}{1.3}\small
\begin{xltabular}{\textwidth}{>{\raggedright\arraybackslash}p{3.2cm}>{\raggedright\arraybackslash}X}
    \caption{Efeito da falha de cada bloco nos fluxos do sistema.}
    \label{tab:fx_falhas}\\
    \hline
    \textbf{Bloco que falha} & \textbf{O que deixa de funcionar e o que continua} \\
    \hline
    \endfirsthead
    \hline
    \textbf{Bloco que falha} & \textbf{O que deixa de funcionar e o que continua} \\
    \hline
    \endhead
    \endfoot
    \multicolumn{2}{c}{Fonte: Autor.} \\
    \endlastfoot
    Front-end & Páginas novas deixam de ser entregues. As telas já abertas continuam a consultar o \emph{back-end}, e a aquisição, o histórico e os alarmes não são afetados. \\
    \hline
    Back-end & Cadastro, login, comando e consultas deixam de responder. A aquisição continua a ler com a última configuração, e o \emph{worker} continua a gravar o valor atual, o histórico e a avaliar os alarmes. \\
    \hline
    \emph{Worker} & O valor atual, o histórico, os alarmes e as variáveis meta deixam de avançar, e as leituras se acumulam nas filas. As telas continuam a responder com o último valor, que passa a desatualizado. \\
    \hline
    Aquisição & Não há leitura nem escrita no equipamento, e os valores passam a desatualizados. O cadastro, as telas e os alarmes já abertos continuam disponíveis, e o comando é recusado. \\
    \hline
    Executor de scripts & Variáveis meta, scripts e componentes de conteúdo param, e a variável meta fica com o último valor. As tags de equipamento, os alarmes sobre elas e os comandos seguem normalmente. \\
    \hline
    Redis & As funções que exigem login e o próprio login respondem indisponibilidade, a aquisição não publica e o comando não executa. Ao voltar vazio, a rotina de reconstrução refaz o valor atual, os estados e a lista de sessões revogadas. \\
    \hline
    PostgreSQL & Cadastro, login, histórico e alarmes deixam de funcionar, e a tela de valores informa indisponibilidade. O valor atual continua a ser atualizado no Redis. \\
    \hline
    Campo & O último valor é mantido e passa a desatualizado, o evento de falha de comunicação do dispositivo é ativado e o comando é recusado. Os demais dispositivos não são afetados. \\
    \hline
\end{xltabular}}

\subsection{Resumo dos Fluxos}

\par A Tabela~\ref{tab:fx_resumo_fluxos} reúne os sete fluxos, os blocos que cada um percorre e o que transporta o trabalho entre eles. A comunicação entre os serviços usa dois mecanismos: a chamada do Moleculer, quando o solicitante precisa da resposta, como no comando, e a fila, quando o trabalho pode esperar e não pode ser perdido, como no histórico. Entre o navegador e o \emph{back-end}, o transporte é o \sigla{http} com dados em \sigla{json}.

{\renewcommand{\arraystretch}{1.3}\small
\begin{xltabular}{\textwidth}{>{\raggedright\arraybackslash}p{3.2cm}>{\raggedright\arraybackslash}X>{\raggedright\arraybackslash}p{4.6cm}}
    \caption{Resumo dos fluxos do sistema.}
    \label{tab:fx_resumo_fluxos}\\
    \hline
    \textbf{Fluxo} & \textbf{Blocos envolvidos} & \textbf{Fila ou ação usada} \\
    \hline
    \endfirsthead
    \hline
    \textbf{Fluxo} & \textbf{Blocos envolvidos} & \textbf{Fila ou ação usada} \\
    \hline
    \endhead
    \endfoot
    \multicolumn{3}{c}{Fonte: Autor.} \\
    \endlastfoot
    Cadastro de uma tag & Navegador, Front-end, Back-end, Redis, PostgreSQL, Aquisição & \sigla{http}, evento de mudança do cadastro e ações de registro do Moleculer \\
    \hline
    Do campo ao banco & Aquisição, Campo, Redis, \emph{Worker}, PostgreSQL & Ação de leitura do Moleculer, filas de ingestão e de histórico, fila geral de tarefas periódicas \\
    \hline
    Valor atual em tempo real & Navegador, Back-end, Redis & \sigla{http} com consulta periódica, leitura direta do Redis \\
    \hline
    Histórico e tendência & Navegador, Back-end, PostgreSQL & \sigla{http}, consulta somente leitura \\
    \hline
    Comando de equipamentos & Navegador, Back-end, Redis, PostgreSQL, Aquisição, Campo & \sigla{http}, ação de escrita do Moleculer, evento de escrita concluída \\
    \hline
    Alarme & Aquisição, Redis, \emph{Worker}, PostgreSQL, Back-end, Navegador & Fila de ingestão, fila geral de tarefas, sumário no Redis \\
    \hline
    Variável meta & Aquisição, Redis, \emph{Worker}, Executor de scripts, PostgreSQL & Fila de ingestão, fila de ondas de cálculo, ação de avaliação do Moleculer \\
    \hline
\end{xltabular}}

\par O que precisa de resposta imediata passa pelo \emph{back-end}, que sempre confere a permissão. O que não pode se perder passa por uma fila e é processado por um \emph{worker}. O valor atual fica no Redis e o histórico no PostgreSQL, de modo que a tela em tempo real nunca consulta o banco. Em cada fluxo, a falha de uma peça leva à recusa do pedido, à repetição do trabalho ou à marcação do valor como desatualizado, sem produzir um dado falso.
